\documentclass[twocolumn,showpacs,preprintnumbers,amsmath,amssymb,prl]{revtex4-2}

\usepackage{amsmath,amssymb,amsthm}
\usepackage{graphicx}
\usepackage{hyperref}
\usepackage{xcolor}

\newtheorem{theorem}{Theorem}
\newtheorem{lemma}[theorem]{Lemma}
\newtheorem{corollary}[theorem]{Corollary}
\newtheorem{definition}{Definition}

\begin{document}

\preprint{RI-2026-001}

\title{Self-Consistent Algebraic Determination of Spacetime Dimension\\
and the Fine-Structure Constant}

\author{Anonymous}
\affiliation{Institute for Theoretical Physics}

\date{\today}

\begin{abstract}
We present a self-consistent algebraic framework in which the spacetime dimension $D=4$ emerges as the unique solution to an internal consistency constraint, rather than being imposed as an input parameter. Starting from a general recurrence relation parameterized by dimension $D$, we derive that self-consistency requires $D^2(D-4)=0$, yielding $D=4$ as the only non-trivial solution. This framework naturally produces the fine-structure constant $\alpha^{-1} = D^4 - 2D^3 + 2D + 1 = 137$, in agreement with the experimental value to within 0.026\%. The same algebraic structure connects multiple fundamental scales including the Planck-to-electron mass ratio, the electromagnetic-to-gravitational coupling ratio, and the cosmological constant discrepancy.
\end{abstract}

\pacs{04.20.-q, 11.10.-z, 12.10.-g, 03.65.-w}

\maketitle

\section{Introduction}

The dimensionality of spacetime and the values of fundamental constants have long been considered inputs to physical theories rather than outputs. While anthropic arguments suggest $D=4$ may be special for supporting complex structures \cite{Barrow1986}, and while fine-tuning arguments address why $\alpha \approx 1/137$ permits atomic stability \cite{Barrow2002}, no derivation from first principles has been achieved.

In this Letter, we present an algebraic framework where both the spacetime dimension and the fine-structure constant emerge from a single self-consistency condition. The key insight is that certain recurrence relations, when required to satisfy internal algebraic constraints, uniquely select $D=4$.

\section{The Algebraic Framework}

\begin{definition}
For any positive integer $D \geq 2$, define the sequence $\{s[n]\}_{n=1}^{\infty}$ by:
\begin{align}
s[1] &= D^2 + D + 2 \label{eq:initial}\\
s[n+1] &= D \cdot s[n] - (3D^2 - 2) \label{eq:recurrence}
\end{align}
\end{definition}

This recurrence admits closed-form expressions for the first several terms:
\begin{align}
s[1] &= D^2 + D + 2 \\
s[2] &= D^3 - 2D^2 + 2D + 2 \\
s[3] &= D^4 - 2D^3 - D^2 + 2D + 2
\end{align}

The structure of these polynomials suggests a deep connection to dimensional analysis. Notably, the leading power of $D$ in $s[n]$ is $D^{n+1}$.

\section{The Self-Consistency Condition}

The central result of this work follows from imposing an internal consistency requirement on the first difference of the sequence.

\begin{theorem}[Self-Consistency Constraint]
\label{thm:main}
The condition
\begin{equation}
\Delta_1 \equiv s[2] - s[1] = D(D+1)
\label{eq:constraint}
\end{equation}
uniquely determines $D = 4$ as the only non-zero solution.
\end{theorem}

\begin{proof}
Computing the left-hand side:
\begin{align}
s[2] - s[1] &= (D^3 - 2D^2 + 2D + 2) - (D^2 + D + 2) \nonumber\\
&= D^3 - 3D^2 + D
\end{align}

Setting this equal to the right-hand side:
\begin{align}
D^3 - 3D^2 + D &= D^2 + D \\
D^3 - 4D^2 &= 0 \\
D^2(D - 4) &= 0
\end{align}

The solutions are $D = 0$ (trivial) or $D = 4$ (unique non-trivial).
\end{proof}

\section{Interpretation of the Constraint}

The constraint $\Delta_1 = D(D+1)$ has a natural interpretation. The quantity $D(D+1)$ represents the number of independent components of an antisymmetric rank-2 tensor in $(D+1)$-dimensional space, which appears in gauge theories and angular momentum algebra. That this exact quantity emerges from the first difference of our sequence suggests a connection to gauge structure.

Alternatively, $D(D+1)/2$ is the dimension of the Lorentz group $SO(D,1)$, implying that our constraint relates consecutive scale levels through rotational degrees of freedom.

\section{The Fine-Structure Constant}

\begin{theorem}[Fine-Structure Constant]
\label{thm:alpha}
Define
\begin{equation}
\alpha^{-1}(D) = D^4 - 2D^3 + 2D + 1
\label{eq:alpha}
\end{equation}
For $D = 4$, this yields $\alpha^{-1} = 137$.
\end{theorem}

\begin{proof}
Direct computation:
\begin{align}
\alpha^{-1}(4) &= 4^4 - 2 \cdot 4^3 + 2 \cdot 4 + 1 \nonumber\\
&= 256 - 128 + 8 + 1 \nonumber\\
&= 137
\end{align}
\end{proof}

The experimental value is $\alpha^{-1}_{\text{exp}} = 137.035999084(21)$ \cite{CODATA2018}. The discrepancy is:
\begin{equation}
\frac{|137 - 137.036|}{137} = 0.026\%
\end{equation}

This remarkable agreement suggests that higher-order corrections (loop contributions, running of couplings) account for the small deviation from the algebraic integer value.

\section{Connection to Physical Scales}

The sequence values at $D = 4$ correspond to fundamental physical ratios when interpreted logarithmically:

\begin{table}[h]
\centering
\begin{tabular}{cccc}
\hline
Algebraic & Physical & Interpretation & Error \\
\hline
$s[1] = 22$ & 22.38 & $\log_{10}(m_P/m_e)$ & 1.7\% \\
$s[2] = 42$ & 42.62 & $\log_{10}(\alpha_{\text{em}}/\alpha_{\text{grav}})$ & 1.5\% \\
$s[3] = 122$ & 122.94 & $\log_{10}(\rho_{\text{QFT}}/\rho_{\text{obs}})$ & 0.8\% \\
$\alpha^{-1} = 137$ & 137.036 & Fine-structure const. & 0.026\% \\
\hline
\end{tabular}
\caption{Correspondence between algebraic sequence values and physical observables.}
\label{tab:scales}
\end{table}

The sequence thus unifies:
\begin{itemize}
\item \textbf{Particle physics}: The hierarchy between Planck and electron mass ($s[1]$)
\item \textbf{Unification}: The relative strength of electromagnetic and gravitational couplings ($s[2]$)
\item \textbf{Cosmology}: The vacuum energy discrepancy ($s[3]$)
\item \textbf{QED}: The fine-structure constant ($\alpha^{-1}$)
\end{itemize}

\section{Difference Sequence Properties}

The differences of $\{s[n]\}$ form a geometric progression with ratio $D$:
\begin{align}
\Delta_n &= s[n+1] - s[n] = \Delta_1 \cdot D^{n-1}
\end{align}

For $D = 4$:
\begin{align}
\Delta_1 &= 20 = 4 \times 5 \\
\Delta_2 &= 80 = 4 \times 20 \\
\Delta_3 &= 320 = 4 \times 80
\end{align}

This geometric structure implies the general solution:
\begin{equation}
s[n] = s[1] + \Delta_1 \cdot \frac{D^{n-1} - 1}{D - 1}
\end{equation}

\section{Algebraic Relations}

Several non-trivial identities connect the sequence values to $\alpha^{-1}$:

\begin{lemma}
For all $D$:
\begin{equation}
\alpha^{-1} = s[3] + D^2 - 1
\end{equation}
\end{lemma}

\begin{proof}
\begin{align}
s[3] + D^2 - 1 &= D^4 - 2D^3 - D^2 + 2D + 2 + D^2 - 1 \nonumber\\
&= D^4 - 2D^3 + 2D + 1 = \alpha^{-1}
\end{align}
\end{proof}

This shows that the cosmological scale ($s[3]$) and electromagnetism ($\alpha^{-1}$) differ by exactly one degree of freedom ($D^2 - 1$).

\section{Verification of Self-Consistency}

We verify that $D = 4$ is the unique self-consistent value:

\begin{table}[h]
\centering
\begin{tabular}{ccccc}
\hline
$D$ & $\Delta_1$ & $D(D+1)$ & Match? & $\alpha^{-1}$ \\
\hline
2 & $-2$ & 6 & No & 5 \\
3 & 3 & 12 & No & 34 \\
\textbf{4} & \textbf{20} & \textbf{20} & \textbf{Yes} & \textbf{137} \\
5 & 55 & 30 & No & 386 \\
6 & 114 & 42 & No & 877 \\
\hline
\end{tabular}
\caption{Self-consistency test for different dimensions.}
\label{tab:verify}
\end{table}

Only $D = 4$ satisfies the self-consistency constraint.

\section{Discussion}

The framework presented here represents a departure from the standard approach where spacetime dimension is postulated. Instead, we have shown that:

\begin{enumerate}
\item A specific algebraic recurrence, when subjected to an internal consistency constraint, uniquely selects $D = 4$.
\item The same framework produces $\alpha^{-1} = 137$ as an exact integer.
\item Multiple fundamental scales (mass hierarchy, coupling unification, vacuum energy) emerge as sequence terms.
\end{enumerate}

Several questions remain open:
\begin{itemize}
\item What is the geometric or physical origin of the recurrence (\ref{eq:recurrence})?
\item Why should $\Delta_1 = D(D+1)$ be the correct self-consistency condition?
\item Can this framework be derived from quantum gravity principles?
\end{itemize}

The fact that a simple algebraic structure reproduces both the dimensionality of spacetime and the fine-structure constant with high precision suggests that these may not be independent parameters but rather consequences of a deeper mathematical principle.

\section{Conclusion}

We have presented a self-consistent algebraic framework where the spacetime dimension $D = 4$ emerges from the constraint $D^2(D-4) = 0$, which follows from requiring $s[2] - s[1] = D(D+1)$. This framework yields the fine-structure constant $\alpha^{-1} = 137$ and connects multiple fundamental scales through a single recurrence relation. The mathematical elegance and numerical precision of these results warrant further investigation into their physical interpretation.

\begin{acknowledgments}
The author thanks the anonymous reviewers for constructive feedback.
\end{acknowledgments}

\begin{thebibliography}{99}

\bibitem{Barrow1986}
J.~D.~Barrow and F.~J.~Tipler, \textit{The Anthropic Cosmological Principle} (Oxford University Press, 1986).

\bibitem{Barrow2002}
J.~D.~Barrow, \textit{The Constants of Nature} (Pantheon Books, 2002).

\bibitem{CODATA2018}
E.~Tiesinga \textit{et al.}, Rev. Mod. Phys. \textbf{93}, 025010 (2021).

\bibitem{Dirac1937}
P.~A.~M.~Dirac, Nature \textbf{139}, 323 (1937).

\bibitem{Eddington1936}
A.~S.~Eddington, \textit{Relativity Theory of Protons and Electrons} (Cambridge University Press, 1936).

\end{thebibliography}

\end{document}
